\documentclass[10pt,a4paper]{article}
\usepackage[utf8]{inputenc}
\usepackage[T1]{fontenc}
\usepackage[ngerman]{babel}
\usepackage[margin=1.7cm,top=1.5cm,bottom=1.7cm]{geometry}
\usepackage{titlesec}
\titleformat*{\section}{\normalfont\large\bfseries}
\titleformat*{\subsection}{\normalfont\normalsize\bfseries}
\titlespacing*{\section}{0pt}{8pt}{3pt}
\titlespacing*{\subsection}{0pt}{6pt}{2pt}
\usepackage{amsmath,amssymb}
\usepackage{listings}
\usepackage{xcolor}
\usepackage{fancyhdr}
\usepackage{parskip}
\usepackage{needspace}
\usepackage{tikz}
\usepackage{array}
\usepackage{enumitem}

\definecolor{loesfarbe}{RGB}{0,110,90}
\definecolor{neurot}{RGB}{200,30,60}

\newcommand{\mcbox}{\raisebox{-1pt}{\fbox{\rule{0pt}{7pt}\rule{7pt}{0pt}}}\hspace{6pt}}
\newcommand{\mcja}{\raisebox{-1pt}{{\fboxrule0.9pt\fcolorbox{loesfarbe}{white}{%
  \rule{0pt}{7pt}\makebox[7pt][c]{\color{loesfarbe}\scriptsize\bfseries X}}}}\hspace{6pt}}
\newcommand{\pkt}[1]{\hfill\textbf{(#1 P)}}
\newcommand{\neu}{{\color{neurot}\footnotesize\bfseries[NEU]}\ }

\newenvironment{loes}%
  {\par\vspace{2pt}\begingroup\small\noindent\textbf{\color{loesfarbe}$\blacktriangleright$ Lösung.}\ }%
  {\endgroup\par\vspace{5pt}}

\lstset{language=Python,basicstyle=\ttfamily\footnotesize,keywordstyle=\bfseries,
  commentstyle=\itshape\color{gray},showstringspaces=false,frame=single,
  framerule=0.3pt,xleftmargin=8pt,framexleftmargin=6pt,aboveskip=6pt,belowskip=6pt}

\pagestyle{fancy}
\fancyhf{}
\fancyhead[L]{\small MDT5/2 -- Probeklausur 2 \textbf{\color{loesfarbe}mit Lösungen}}
\fancyhead[R]{\small Maschinelles Lernen zur Anomalieerkennung}
\fancyfoot[C]{\small Seite \thepage}
\renewcommand{\headrulewidth}{0.4pt}

\setlist[itemize]{leftmargin=12pt,itemsep=1pt,topsep=2pt}
\setlist[enumerate]{leftmargin=16pt,itemsep=1pt,topsep=2pt}

\begin{document}

\begin{center}
  {\Large\bfseries Probeklausur 2 -- ausgefüllt}\\[4pt]
  {\large Maschinelles Lernen zur Anomalieerkennung -- Kapitel 1--8}\\[6pt]
  \begin{tabular}{@{}ll@{}}
  \textbf{Gesamtpunktzahl:} & 125 Punkte\\
  \end{tabular}
\end{center}

\vspace{2pt}
\noindent\fcolorbox{neurot}{neurot!5}{\parbox{\dimexpr\linewidth-2\fboxsep-2\fboxrule}{%
\textbf{\color{neurot}Lesehilfe.} \mcja{} markiert die richtige Aussage.
Mit \neu sind die acht Stellen markiert, die \emph{nicht} schon auf deinen Blättern stehen:
Dense auf 4D wirft keinen Fehler (3.2) $\cdot$ AUC-Vorzeichen (3.1) $\cdot$ $\pm1$ vs.\ $0/1$
(3.1) $\cdot$ Alarmzahl gegenrechnen (5c) $\cdot$ Preis je Ausreißer (5d) $\cdot$ Punkt exakt auf
dem Margin-Rand (6.1c) $\cdot$ linearer Kernel übersieht Seitliches (6.3a) $\cdot$ Elliptic
Envelope hält $P_1$ für den normalsten Punkt (7b).}}

\noindent\rule{\linewidth}{0.8pt}

\section*{Aufgabe 1 -- Multiple Choice mit Begründung \pkt{21}}

\subsection*{1.1 \ Für die Normalisierung von Merkmalen gilt:}
\mcbox Eine z-Transformation (\texttt{StandardScaler}) bildet die Merkmale \textbf{immer} auf den
Bereich $[-1, 1]$ ab.

\mcbox Für den Isolation Forest ist eine Normalisierung \textbf{essentiell}, weil er die Abstände
zwischen den Punkten auswertet.

\mcja Für abstandsbasierte Verfahren ist eine Normalisierung nötig, weil sonst das Merkmal mit
dem größten Wertebereich den Abstand dominiert.

\begin{loes}
In jeden Abstand gehen alle Merkmale mit ihrer \emph{rohen} Größenordnung ein; ein Merkmal in
Kilogramm überstimmt eines in Metern allein durch seinen Zahlenbereich.
Aussage 1 ist falsch: die z-Transformation setzt Mittelwert $0$ und Standardabweichung $1$ und hat
\textbf{keinen festen Wertebereich} ($[0,1]$ bzw.\ $[-1,1]$ liefert die Min-Max-Normalisierung).
Aussage 2 ist falsch: der Isolation Forest wertet \textbf{keine Abstände} aus, sondern splittet je
ein Merkmal einzeln -- Normalisierung ist dort überflüssig.
\end{loes}

\subsection*{1.2 \ Für den Parameter \texttt{contamination} der Elliptic Envelope gilt:}
\mcbox \texttt{contamination} gibt die Anzahl der erwarteten Ausreißer im Datensatz an und
verändert den Kovarianz-Fit.

\mcja \texttt{contamination} ist ein Anteil und wirkt als Schwellwert; die AUC des Modells
ändert sich dadurch \textbf{nicht}.

\mcbox \texttt{contamination} ist ein Formparameter der geschätzten Ellipse und lässt sich
unüberwacht per Grid Search bestimmen.

\begin{loes}
\texttt{score\_samples} (der negative Mahalanobis-Abstand) hängt \textbf{nur vom Kovarianz-Fit} ab;
\texttt{contamination} verschiebt allein \texttt{offset\_}, also die $\pm1$-Grenze in
\texttt{predict}. Die Scores selbst ändern sich nicht, und die AUC ist ohnehin schwellwertfrei --
sie bleibt doppelt unberührt. Es ist ein \textbf{Anteil}, keine Anzahl (Aussage 1), und kein
Formparameter (Aussage 3); unüberwacht ist er nicht schätzbar, weil die Suchkurven flach bleiben.
Merkbild: AUC $=$ feste Modellgüte, \texttt{contamination} $=$ Betriebspunkt auf der festen
ROC-Kurve -- dasselbe Muster wie $\nu$ bei Deep SVDD.
\end{loes}

\subsection*{1.3 \ Für einen Autoencoder zur Anomalieerkennung gilt:}
\mcbox Die Sigmoid-Aktivierung in der Ausgabeschicht zeigt an, dass der Autoencoder ein binäres
Klassifikationsproblem (normal / anomal) löst.

\mcja Hat die Code-Schicht $z$ mindestens so viele Dimensionen wie die Eingabe, kann der
Autoencoder die Identität lernen und wird zur Anomalieerkennung unbrauchbar.

\mcbox Der Autoencoder wird auf Normaldaten \textbf{und} Anomalien trainiert, damit er beide
Klassen unterscheiden kann.

\begin{loes}
Der Bottleneck erzwingt Kompression: nur wenn $z$ \textbf{kleiner} als die Eingabe ist, muss das
Netz das Wesentliche der Normaldaten lernen. Reicht die Kapazität für eine Eins-zu-eins-Abbildung,
rekonstruiert der AE \emph{alles} gut, auch Anomalien $\Rightarrow$ der Rekonstruktionsfehler
trennt nicht mehr.
Aussage 1 ist falsch: der AE \textbf{klassifiziert nichts}, er liefert nur einen
Rekonstruktionsfehler; sigmoid steht dort, weil die Pixel auf $[0,1]$ normiert sind.
Aussage 3 ist falsch: trainiert wird \textbf{nur auf Normaldaten}, gerade damit Anomalien schlecht
rekonstruiert werden.
\end{loes}

\subsection*{1.4 \ Für Gradienten in tiefen Netzen gilt:}
\mcbox ReLU verhindert sowohl verschwindende als auch explodierende Gradienten, weil ihre
Ableitung \textbf{nie} größer als 1 wird.

\mcja Die Ableitung der Sigmoidfunktion ist höchstens $0{,}25$; über viele Schichten hinweg
verschwindet der Gradient deshalb in Richtung der \textbf{Eingangsschicht}.

\mcbox Explodierende Gradienten treten \textbf{immer} zuerst in der Ausgangsschicht auf, weil
dort der Fehler entsteht.

\begin{loes}
Backpropagation multipliziert pro Schicht einen Faktor $\sigma'(z) \le 0{,}25$ (Maximum bei
$z=0$); über $L$ Schichten also $0{,}25^L$ -- die vorderen Schichten bekommen praktisch keinen
Gradienten mehr.
Aussage 1: Die Ableitung von ReLU ist tatsächlich nie $>1$ (sie ist exakt $0$ oder $1$), aber
daraus folgt \textbf{nicht}, dass Exploding verhindert wird -- beim Backprop werden zusätzlich
die \textbf{Gewichte} mitmultipliziert. ReLU löst nur das Vanishing-Problem.
Aussage 3: Backprop läuft von hinten nach vorn, also eskalieren \textbf{beide} Phänomene Richtung
\textbf{Eingangsschicht}, nicht zur Ausgangsschicht.
\end{loes}

\subsection*{1.5 \ Für die One-Class SVM gilt:}
\mcbox $\nu$ ist ein Strafgewicht für Fehlklassifikationen (wie $C$ bei der Zweiklassen-SVM):
je größer $\nu$, desto härter die Grenze.

\mcja $\nu$ ist eine obere Schranke für den Anteil der Ausreißer im Training und zugleich eine
untere Schranke für den Anteil der Support-Vektoren.

\mcbox Der Kernel-Trick ist bei der One-Class SVM nur eine Beschleunigung -- man könnte die
Transformation $\varphi(x)$ \textbf{immer} auch explizit ausrechnen.

\begin{loes}
$\nu$ ist eine \textbf{Mengenvorgabe}, kein Strafgewicht -- und es sind \textbf{zwei} Schranken:
höchstens der Anteil $\nu$ der Trainingspunkte darf außerhalb liegen, mindestens der Anteil $\nu$
sind Support-Vektoren. Deshalb sucht man $\nu \in (0;1]$ auch \textbf{linear} (\texttt{linspace}),
nicht logarithmisch.
Aussage 1 dreht zusätzlich die Richtung um: $C$ und $\nu$ laufen \textbf{gegenläufig} -- großes $C$
entspricht \textbf{kleinem} $\nu$.
Aussage 3 ist falsch: beim RBF-Kernel hat $\varphi(x)$ \textbf{unendlich viele} Komponenten, es
ließe sich gar nicht ausrechnen -- der Kernel-Trick ist dort die einzige Möglichkeit, nicht bloß
die schnellere.
\end{loes}

\subsection*{1.6 \ Ein Detektor flaggt in einem Datensatz mit 2\,\% Anomalien sehr viele Punkte.
Dazu gilt:}
\mcbox Die \emph{balanced accuracy} ist hier die ehrlichste Metrik, weil sie die Unbalanciertheit
der Klassen einbezieht.

\mcja Precision und $F_1$ decken die Fehlalarme auf, weil in ihnen \textbf{keine} True Negatives
vorkommen.

\mcbox Ein hoher Recall belegt, dass nur wenige Fehlalarme auftreten.

\begin{loes}
Die ehrlichste Metrik ist nicht die, die Unbalanciertheit \emph{einbezieht}, sondern die, die die
\textbf{große Klasse gar nicht enthält}. Bei 2\,\% Anomalien sind 98\,\% der Punkte normal; die
vielen TN blähen Accuracy auf und beschönigen auch die balanced accuracy noch. Precision
($\text{TP}/(\text{TP}+\text{FP})$) und $F_1$ kennen \textbf{keine TN} -- massenhafte Fehlalarme
schlagen dort sofort durch.
Aussage 3 ist der klassische Tausch: \textbf{Recall} $=$ von den echten Anomalien, wie viele
gefangen. Er kann $1{,}0$ sein, während der Detektor \emph{alles} flaggt. „Von den geflaggten, wie
viele echt`` ist die \textbf{Precision}.
\end{loes}

\subsection*{1.7 \ Für Deep SVDD und GOAD gilt:}
\mcbox Batch Normalization ist im Deep-SVDD-Netz \textbf{grundsätzlich verboten}, weil der
gelernte Parameter $\beta$ wie ein Bias wirkt.

\mcja $\nu$ legt bei Deep SVDD den Radius als $(1-\nu)$-Quantil der Abstände zu $c$ fest; es
verschiebt nur den Betriebspunkt, die AUC bleibt gleich.

\mcbox Der dritte Term der GOAD-Verlustfunktion ist ein Weight Decay auf die Gewichte des Netzes.

\begin{loes}
$\nu$ kommt im Training gar nicht vor -- die vereinfachte Zielfunktion minimiert nur den Abstand zu
$c$. Erst \textbf{danach} legt $\nu$ den Radius als $(1-\nu)$-Quantil der Trainingsabstände fest,
verschiebt also nur die $\pm1$-Grenze auf einer bereits feststehenden ROC-Kurve. Größeres $\nu$
$\Rightarrow$ kleinerer Radius $\Rightarrow$ mehr Alarme. Exakt dasselbe Muster wie
\texttt{contamination} in 1.2.
Aussage 1 ist falsch: BatchNorm ist \textbf{erlaubt} und steht in den Musterlösungen wie im
Praktikumscode. Verboten sind Biases und gedeckelte Aktivierungen.
Aussage 3 ist falsch: der dritte GOAD-Term ist eine \textbf{$L_2$-Regularisierung der
Latent-Vektoren $z$}, kein Weight Decay auf die Gewichte.
\end{loes}

\newpage
\section*{Aufgabe 2 -- Ausgabe-Shapes eines Faltungs-Autoencoders \pkt{12}}

\begin{lstlisting}
ae = tf.keras.Sequential([
    tf.keras.layers.Input((64, 64, 3)),                                   #  ( 64, 64, 3 )
 1  tf.keras.layers.Conv2D(16, (3,3), padding='same', activation='relu'),
 2  tf.keras.layers.MaxPooling2D((2,2)),
 3  tf.keras.layers.Conv2D(32, (5,5), padding='valid', activation='relu'),
 4  tf.keras.layers.Conv2D(32, (3,3), strides=2, padding='same', activation='relu'),
 5  tf.keras.layers.MaxPooling2D((3,3)),
 6  tf.keras.layers.Flatten(),
 7  tf.keras.layers.Dense(20, activation='relu'),
 8  tf.keras.layers.Dense(512, activation='relu'),
 9  tf.keras.layers.Reshape((4, 4, 32)),
10  tf.keras.layers.Conv2DTranspose(32, (3,3), strides=2, padding='same', activation='relu'),
11  tf.keras.layers.Conv2DTranspose(16, (3,3), strides=2, padding='same', activation='relu'),
12  tf.keras.layers.Conv2DTranspose(8,  (3,3), strides=2, padding='same', activation='relu'),
13  tf.keras.layers.Conv2DTranspose(3,  (3,3), strides=2, padding='same', activation='sigmoid'),
])
\end{lstlisting}

\begin{center}
\begin{tabular}{|c|p{4.2cm}||c|p{4.2cm}|}
\hline
Zeile & Ausgabe-Shape & Zeile & Ausgabe-Shape \\ \hline
1 & \rule{0pt}{13pt}\color{loesfarbe}$64\times64\times16$ & 8 & \rule{0pt}{13pt}\color{loesfarbe}$512$\\ \hline
2 & \rule{0pt}{13pt}\color{loesfarbe}$32\times32\times16$ & 9 & \rule{0pt}{13pt}\color{loesfarbe}$4\times4\times32$\\ \hline
3 & \rule{0pt}{13pt}\color{loesfarbe}$28\times28\times32$ & 10 & \rule{0pt}{13pt}\color{loesfarbe}$8\times8\times32$\\ \hline
4 & \rule{0pt}{13pt}\color{loesfarbe}$14\times14\times32$ & 11 & \rule{0pt}{13pt}\color{loesfarbe}$16\times16\times16$\\ \hline
5 & \rule{0pt}{13pt}\color{loesfarbe}$4\times4\times32$ & 12 & \rule{0pt}{13pt}\color{loesfarbe}$32\times32\times8$\\ \hline
6 & \rule{0pt}{13pt}\color{loesfarbe}$512$ & 13 & \rule{0pt}{13pt}\color{loesfarbe}$64\times64\times3$\\ \hline
7 & \rule{0pt}{13pt}\color{loesfarbe}$20$ & & \\ \hline
\end{tabular}
\end{center}

\begin{loes}
Die drei Regeln: \texttt{same} mit strides 1 $\Rightarrow$ Größe \textbf{bleibt};
\texttt{valid} mit Kernel $k$ $\Rightarrow$ \textbf{minus $(k-1)$}; Halbieren $\Rightarrow$
Pooling \textbf{abrunden}, Conv mit \texttt{strides=2} \textbf{aufrunden}.
Zeile 3: $32-4 = 28$. Zeile 4: $\lceil 28/2\rceil = 14$. Zeile 5: $\lfloor 14/3 \rfloor = 4$.\\[2pt]
\textbf{Gegenprobe, die jeden Kettenfehler aufdeckt:} Zeile 9 ist
\texttt{Reshape((4,4,32))} $= 512$ und Zeile 8 ist \texttt{Dense(512)}. Ein Autoencoder ist
symmetrisch, also \textbf{muss} der Flatten in Zeile 6 ebenfalls $512$ ergeben. Kommt dort etwas
anderes heraus, steckt weiter oben ein Rechenfehler -- nicht weiterrechnen, zurückgehen.
\end{loes}

\textbf{b)} Wie viele trainierbare Parameter hat Zeile 1, wie viele Zeile 7? Rechenweg angeben.

\begin{loes}
\textbf{Conv:} $k_h \cdot k_w \cdot C_{\text{in}} \cdot f + f$ \quad --- \quad
\textbf{Dense:} $n_{\text{in}} \cdot n_{\text{out}} + n_{\text{out}}$
\[
\text{Zeile 1:}\quad 3\cdot3\cdot\underbrace{3}_{\text{RGB}}\cdot16 + 16 = 432+16 = \mathbf{448}
\qquad
\text{Zeile 7:}\quad \underbrace{512}_{\text{nach Flatten}}\cdot20 + 20 = \mathbf{10\,260}
\]
$C_{\text{in}}$ sind die \textbf{Kanäle des Eingangstensors} (hier 3), niemals die Bildkantenlänge;
$n_{\text{in}}$ ist die \textbf{Länge nach dem Flatten}, nicht $H$ oder $C$ allein.
\end{loes}

\textbf{c)} Welche Schicht des Netzes trägt das größte Overfitting-Risiko und warum?

\begin{loes}
Die \textbf{vollverknüpften Dense-Schichten} (Zeilen 7/8). Ihre Gewichtszahl ist das
\textbf{Produkt} aus Eingangslänge und Neuronenzahl und wächst mit der Auflösung; Faltungen
brauchen dank \textbf{Weight Sharing} nur $k\cdot k\cdot C_{\text{in}}\cdot C_{\text{out}}$
\emph{unabhängig} von der Bildgröße. Viele Parameter $\Rightarrow$ hohe Overfitting-Gefahr,
deshalb sitzen Dropout und $L_2$ genau dort.
Nicht „weil sie die Identität lernen könnte`` -- das ist das Bottleneck-Argument aus 1.3 und
beantwortet eine andere Frage.
\end{loes}

\textbf{d)} Begründen Sie \textbf{beide} Entwurfsentscheidungen in Zeile 13: warum genau 3 Filter
und warum \texttt{sigmoid}?

\begin{loes}
\textbf{3 Filter,} weil die letzte transponierte Faltung so viele Feature Maps haben muss, wie das
Zielbild \textbf{Kanäle} hat -- RGB, also 3 (bei MNIST wäre es 1). Sonst passt die Ausgabe-Shape
nicht auf die Eingabe, und der Rekonstruktionsfehler ließe sich gar nicht bilden.
\textbf{Sigmoid,} weil die Pixel auf $[0,1]$ normiert sind und sigmoid den Ausgabebereich genau
darauf begrenzt. Das hat \textbf{nichts} mit Klassifikation zu tun.
\end{loes}

\newpage
\section*{Aufgabe 3 -- Fehler finden \pkt{18}}

\subsection*{3.1 \ Training und Evaluierung einer One-Class SVM (Novelty Detection) \pkt{10}}

\begin{lstlisting}
X, y = load_sensor_data()        # y: 0 = normal, 1 = Anomalie; ca. 3 % Anomalien

scaler = StandardScaler()                                                    # (1)
X = scaler.fit_transform(X)                                                  # (2)

X_train, X_test, y_train, y_test = train_test_split(X, y, test_size=0.3)     # (3)

param_grid = {'nu':    np.logspace(-3, 0, 10),                               # (4)
              'gamma': np.logspace(-2, 2, 10)}                               # (5)
search = GridSearchCV(OneClassSVM(kernel='rbf'), param_grid,
                      scoring='accuracy', cv=5)                              # (6)
search.fit(X_train, y_train)                                                 # (7)

best = search.best_estimator_
pred = best.predict(X_train)                                                 # (8)
pred = np.where(pred == 1, 0, 1)
print('Accuracy:', accuracy_score(y_train, pred))                            # (9)
print('AUC:', roc_auc_score(y_test, best.score_samples(X_test)))             # (10)
\end{lstlisting}

\begin{loes}
\begin{enumerate}
\item \textbf{Zeile (2) -- Data Leakage.} Der Scaler wird auf den \textbf{Gesamtdaten} gefittet,
  \emph{bevor} gesplittet wird. Mittelwert und Streuung der Testdaten fließen ins Training ein.
  Richtig: erst splitten, dann \texttt{fit} nur auf \texttt{X\_train},
  \texttt{transform} auf beide.
\item \textbf{Zeile (7) -- Novelty-Verfahren auf kontaminierten Daten.} Die OCSVM muss auf
  \textbf{reinen Normaldaten} gefittet werden; \texttt{X\_train} enthält aber die 3\,\% Anomalien.
  So wird aus Novelty unbeabsichtigt \textbf{Outlier} Detection, die Grenze wird zu den Anomalien
  hin aufgeweitet, die TPR bricht ein. Richtig: nur \texttt{X\_train[y\_train == 0]} fitten.
  Zusätzlich wird \texttt{y\_train} an ein \textbf{unüberwachtes} \texttt{fit} übergeben.
\item \textbf{Zeile (6) -- kein Pipeline-Objekt.} Ohne
  \texttt{Pipeline(StandardScaler(), OneClassSVM())} wiederholt sich die Leakage aus (2) in
  \textbf{jedem CV-Fold}: der Validierungsteil ist bereits mitskaliert.
\item \textbf{Zeile (6) -- \texttt{scoring='accuracy'}.} Bei 3\,\% Anomalien erreicht „alles
  normal`` schon 97\,\%; die Suche optimiert also auf eine Metrik, die nichts unterscheidet.
  Besser \texttt{balanced\_accuracy} oder AUC.
  \neu Zusätzlich passen die \textbf{Kodierungen} nicht: \texttt{predict} liefert $\pm1$,
  \texttt{y\_train} ist $0/1$ -- verglichen werden zwei verschiedene Skalen.
\item \textbf{Zeilen (8)/(9) -- Bewertung auf Trainingsdaten.} \texttt{predict} und
  \texttt{accuracy\_score} laufen auf \texttt{X\_train}/\texttt{y\_train}, also auf genau den
  Daten, aus denen das Modell \emph{und} die Metaparameter stammen. Muss \texttt{X\_test} sein.
\item \neu \textbf{Zeile (10) -- Vorzeichen.} \texttt{score\_samples} ist \textbf{hoch = normal},
  \texttt{y\_test} ist aber \textbf{1 = Anomalie}. Die AUC kommt gespiegelt heraus (unter 0,5).
  Richtig: \texttt{roc\_auc\_score(y\_test, -best.score\_samples(X\_test))}.
\item \textbf{Zeile (4) -- Raster für $\nu$.} $\nu \in (0;1]$ ist eine Mengenvorgabe, also
  \textbf{linear} rastern (\texttt{np.linspace}); $10^0 = 1$ ist zudem der entartete Grenzfall.
\item \emph{Nebenbei:} kein \texttt{random\_state} (und kein \texttt{stratify}) im Split
  $\Rightarrow$ nicht reproduzierbar.
\end{enumerate}
\textbf{Kein Fehler} ist \texttt{np.where(pred == 1, 0, 1)} -- das ist die korrekte Umrechnung von
der sklearn-Konvention ($+1$ normal) auf die Label-Konvention ($1 =$ Anomalie).
\end{loes}

\newpage
\subsection*{3.2 \ Aufbau eines CNN-Klassifikators für $32\times32\times3$-Bilder, 5 Klassen \pkt{8}}

\begin{lstlisting}
net = tf.keras.Sequential([
    tf.keras.layers.Input((32, 32, 3)),                            # 32 x 32 x  3
    tf.keras.layers.Conv2D(16, (3,3), padding='same', ...),        # 32 x 32 x 16
    tf.keras.layers.MaxPooling2D((2,2)),                           # 16 x 16 x 16
    tf.keras.layers.Conv2D(32, (3,3), padding='same', ...),        # 16 x 16 x 32
    tf.keras.layers.MaxPooling2D((2,2)),                           #  8 x  8 x 32
    tf.keras.layers.Conv2D(64, (3,3), padding='same', ...),        #  8 x  8 x 64
    tf.keras.layers.MaxPooling2D((2,2)),                           #  4 x  4 x 64
    tf.keras.layers.MaxPooling2D((4,4)),                           #  1 x  1 x 64
    tf.keras.layers.Dense(64, activation='relu'),
    tf.keras.layers.Dense(5,  activation='relu'),
])
net.compile(optimizer='adam', loss='binary_crossentropy', metrics=['accuracy'])
net.fit(x_train, y_train, epochs=20, validation_data=(x_test, y_test))
\end{lstlisting}

\begin{loes}
\begin{enumerate}
\item \neu \textbf{\texttt{Flatten()} fehlt} vor \texttt{Dense(64)}. Nach dem letzten Pooling liegt
  ein $1\!\times\!1\!\times\!64$-Tensor an. \texttt{Dense} wirkt nur auf die \textbf{letzte Achse},
  die Ausgabe wird also $1\!\times\!1\!\times\!5$ statt $(5,)$ und passt nicht zu den Labels.
  \textbf{Es gibt keinen Fehler beim Bauen} -- genau deshalb rutscht das durch.
\item \textbf{Ausgabeschicht mit \texttt{'relu'}} statt \textbf{softmax}. ReLU liefert beliebige
  nichtnegative Werte, die sich nicht zu 1 summieren -- keine Wahrscheinlichkeitsverteilung; dazu
  Gradient 0 für alle abgeschnittenen Neuronen.
\item \textbf{\texttt{loss='binary\_crossentropy'} bei 5 Klassen.} Ab drei Klassen muss es
  \texttt{categorical\_crossentropy} heißen (One-Hot-Labels) -- bei Integer-Labels
  entsprechend \texttt{sparse\_}\discretionary{}{}{}\texttt{categorical\_crossentropy}.
\item \textbf{Doppeltes Pooling.} $4\!\times\!4$ mit Pool $(4,4)$ geht zwar exakt auf, wirft aber
  die \textbf{gesamte räumliche Information} weg -- übrig bleibt ein einzelner Wert je Feature Map.
  Ein Pooling-Schritt gehört gestrichen.
\item \textbf{\texttt{validation\_data=(x\_test, y\_test)}.} Die Testdaten werden zur
  Validierung benutzt; nach 20 Epochen ist die Modellwahl auf ihnen erfolgt und der Testwert nicht
  mehr unabhängig.
\end{enumerate}
\end{loes}

\newpage
\section*{Aufgabe 4 -- Programmieren \pkt{16}}

\subsection*{4.1 \ Deep-SVDD-Netz in Keras \pkt{10}}

\begin{lstlisting}
lam = 1e-6
reg = tf.keras.regularizers.l2(lam)

deep_svdd = tf.keras.Sequential([
    tf.keras.layers.Input((28, 28, 1)),                                       # 28 x 28 x  1
    tf.keras.layers.Conv2D(8, (5,5), padding='same', use_bias=False,
                           kernel_regularizer=reg),                           # 28 x 28 x  8
    tf.keras.layers.BatchNormalization(),
    tf.keras.layers.LeakyReLU(),
    tf.keras.layers.MaxPooling2D((2,2)),                                      # 14 x 14 x  8
    tf.keras.layers.Conv2D(16, (5,5), padding='same', use_bias=False,
                           kernel_regularizer=reg),                           # 14 x 14 x 16
    tf.keras.layers.BatchNormalization(),
    tf.keras.layers.LeakyReLU(),
    tf.keras.layers.MaxPooling2D((2,2)),                                      #  7 x  7 x 16
    tf.keras.layers.Flatten(),                                                #       784
    tf.keras.layers.Dense(32, use_bias=False, kernel_regularizer=reg),        #        32
])
\end{lstlisting}

\begin{loes}
\textbf{Die Konstruktionsregeln, an denen die Punkte hängen:}
\begin{itemize}
\item \textbf{\texttt{use\_bias=False} in jeder Schicht} -- der Keras-Default ist \texttt{True}.
  Ein Bias könnte den konstanten Offset nachbauen, mit dem das Netz alles auf $c$ abbildet.
\item \textbf{Keine gedeckelten Aktivierungen} (kein sigmoid, kein tanh) $\Rightarrow$
  \textbf{LeakyReLU}. Gesättigte Aktivierungen erlauben denselben Kollaps.
\item \textbf{Keine Aktivierung auf der letzten Dense} -- sie ist die Ausgabe $\varphi(x)$.
\item $L_2$ mit $\lambda = 10^{-6}$ über \texttt{kernel\_regularizer} in \emph{jeder} Schicht.
\item \textbf{BatchNorm ist erlaubt} und steht so auch in den Musterlösungen -- nicht anstreichen.
\end{itemize}
Alle drei Verbote laufen auf dieselbe Katastrophe hinaus: \textbf{Nullgewichte}, also die konstante
Abbildung.
\end{loes}

\textbf{b)} Schreiben Sie die Zeile, mit der das Zentrum $c$ bestimmt wird. An welcher Stelle des
Programms muss sie stehen, und was passiert, wenn man sie in die Trainingsschleife setzt oder
$c = \mathbf{0}$ wählt?

\begin{lstlisting}
center = get_center(deep_svdd, train_dataset)   # EINMAL, vor der Trainingsschleife
\end{lstlisting}

\begin{loes}
\textbf{Stelle:} nach der Initialisierung bzw.\ dem Vortraining, \textbf{einmal} und
\textbf{außerhalb} der Trainingsschleife. $c$ ist der Mittelwert der Netzausgaben über die
Trainingsdaten und bleibt danach \textbf{fest}.

\textbf{In der Trainingsschleife:} $c$ wandert bei jedem Schritt mit den Ausgaben mit. Dann lässt
sich der Loss trivial minimieren, indem das Netz alles auf einen beliebigen Punkt zieht und $c$
hinterherwandert $\Rightarrow$ \textbf{Kollaps zur konstanten Abbildung}, Loss $\to 0$,
Scores ununterscheidbar, \textbf{AUC $\approx$ 0,5}.

\textbf{$c = \mathbf{0}$:} Dann sind \textbf{Nullgewichte direkt optimal} -- $\varphi(x) = 0$ für
alle $x$ erzeugt Abstand 0. Der Loss ist minimal, das Modell wertlos: Normaldaten und Anomalien
landen auf demselben Punkt.
\end{loes}

\newpage
\subsection*{4.2 \ Anomalie-Score und Schwellwert in NumPy \pkt{6}}

\begin{lstlisting}
X_rec = ae.predict(X_test)

# 1) Rekonstruktionsfehler je Sample
err = np.mean((X_test - X_rec)**2, axis=(1, 2, 3))      # Shape (N,)

# 2) Schwellwert aus dem bekannten Ausreisseranteil: 2 % oben abschneiden
thr = np.percentile(err, 98)

# 3) sklearn-Konvention: +1 = normal, -1 = Anomalie
result = np.where(err > thr, -1, 1)

# 4) Anzahl der Alarme
print('Alarme:', np.sum(result == -1))
\end{lstlisting}

\textbf{b)} Über welche Achsen wird gemittelt und warum? Welches Perzentil und warum?

\begin{loes}
\textbf{Achsen 1, 2, 3} (Höhe, Breite, Kanal) -- alles außer der \textbf{Batch-Achse 0}, damit pro
Sample \textbf{genau eine Zahl} übrig bleibt. Der Score ist eine Größe je Messung, nicht je Pixel.

\textbf{98.\ Perzentil}, weil aus dem Prozess 2\,\% Ausreißer bekannt sind: die obersten 2\,\% der
Scores werden geflaggt.
\emph{Der Zusatz, der die Antwort vollständig macht:} Das garantiert die richtige \textbf{Anzahl}
an Alarmen ($0{,}02\cdot N$), nicht die richtigen \textbf{Punkte} -- es unterstellt, dass die
Ausreißer tatsächlich die höchsten Scores haben.
\end{loes}

\newpage
\section*{Aufgabe 5 -- Auswertung eines Score-Histogramms \pkt{14}}

\begin{center}\small
\begin{tabular}{|l|r|r|r|r|r|r|r|r|r|r|r|r|}
\hline
Klasse & 0--0,5 & 0,5--1 & 1--1,5 & 1,5--2 & 2--2,5 & 2,5--3 & 3--3,5 & 3,5--4 & 4--4,5 & 4,5--5 & 5--5,5 & 5,5--6\\ \hline
Anzahl & 900 & 1500 & 1200 & 700 & 350 & 180 & 90 & 45 & 20 & 10 & 4 & 1\\ \hline
\end{tabular}
\hspace{4pt}
\begin{tabular}{|c|c|c|c|c|c|}
\hline
$s_0$ & 1,0 & 2,0 & 2,5 & 3,5 & 4,0\\ \hline
TPR & 0,99 & 0,92 & 0,86 & 0,21 & 0,06\\ \hline
FPR & 0,62 & 0,14 & 0,05 & 0,013 & 0,007\\ \hline
\end{tabular}
\quad AUC $=0{,}95$, \ $N = 5000$
\end{center}

\textbf{a)} Schwellwert nur aus dem Histogramm, Anzahl der Alarme, Plausibilität. \pkt{4}

\begin{loes}
Die Verteilung gipfelt bei $s\approx0{,}75$ und fällt monoton. Man legt den Schwellwert dort, wo
sie in den \textbf{flachen Schwanz} übergeht: $\mathbf{s_0 = 3{,}0}$.
\[
90+45+20+10+4+1 = \mathbf{170} \text{ Alarme} \quad\Longrightarrow\quad 170/5000 = 3{,}4\,\%
\]
Plausibel für eine Ausreißerquote. Zum Vergleich: $s_0=2{,}5$ gäbe 350 Alarme (7\,\% -- zu viel),
$s_0 = 3{,}5$ gäbe 80 (1,6\,\%).
\emph{Wichtig als Zusatz:} Aus dem Histogramm \textbf{sieht} man den Übergang gar nicht -- die
entscheidende Region ist genau die, in der die Balken schon auf der Nulllinie liegen (4 bzw.\ 1
Punkt bei einer $y$-Achse bis 1500). Jede Wahl hier ist eine Schätzung.
\end{loes}

\textbf{b)} 98,4\,\%-Quantil: welcher Schwellwert, wie viele Punkte, was ist garantiert? \pkt{4}

\begin{loes}
$1{,}6\,\%$ von 5000 $=$ \textbf{80 Punkte}. Von oben aufsummiert: $1+4+10+20+45 = 80$ -- das sind
exakt die Klassen ab 3,5, also $\mathbf{s_0 = 3{,}5}$.

\textbf{Garantiert} ist die \textbf{Anzahl} der Alarme (genau 1,6\,\%, die Alarmlast ist planbar).
\textbf{Nicht garantiert} ist, dass es die \textbf{richtigen} Punkte sind: das Quantil unterstellt
stillschweigend, die Ausreißer seien die 80 höchsten Scores -- geprüft wird das nie. Die
Betriebspunkt-Tabelle zeigt das Gegenteil: bei 3,5 ist TPR nur $0{,}21$.
\end{loes}

\textbf{c)} Nur 21\,\% gefunden bei AUC 0,95 -- Modell oder Schwellwert? \pkt{4}

\begin{loes}
\textbf{Der Schwellwert.} Die AUC ist schwellwertfrei und bewertet allein die \textbf{Rangfolge}
der Scores -- die ist mit 0,95 gut. Ein einzelner Betriebspunkt sagt nichts über die Modellgüte.

Verantwortlich ist der \textbf{schwere obere Schwanz der Normaldaten}: die Normalpunkte reichen mit
einem langen rechten Ausläufer durch das Band der Ausreißer hindurch und darüber hinaus. Damit
bestehen die obersten 1,6\,\% der Scores \textbf{überwiegend aus Normaldaten}, und das Quantil
schneidet oberhalb des Anomalie-Hauptfelds ab.

\neu \textbf{Gegenrechnung -- TP $+$ FP muss die geflaggte Menge treffen} (80 Anomalien,
4920 Normale):
\[
\text{TP} = 0{,}21\cdot80 \approx 17,\qquad
\text{FP} = 0{,}013\cdot4920 \approx 64,\qquad
\text{TP}+\text{FP} \approx 81 \approx 80 \ \checkmark
\]
Von den 80 Alarmen sind rund \textbf{drei Viertel Fehlalarme}. Stimmt die Summe nicht, steckt der
Fehler in der Quantilrechnung -- eine direkte Selbstkontrolle in der Klausur.

\emph{Diagnoseregel:} hohe AUC $+$ schlechte TPR $\Rightarrow$ falsche \textbf{Schwelle};
niedrige AUC $\Rightarrow$ schlechtes \textbf{Modell}.
\end{loes}

\textbf{d)} Bestes $s_0$ im Sinne des ROC-Knies; was kostet der letzte Rest TPR? \pkt{5}

\begin{loes}
Rechenkriterium: kleinster Abstand zur idealen Ecke $(0,1)$, also Minimum von
$d^2 = \text{FPR}^2 + (1-\text{TPR})^2$.
\begin{center}\small
\begin{tabular}{|c|c|c|c|c|c|}
\hline
$s_0$ & 1,0 & 2,0 & \textbf{2,5} & 3,5 & 4,0\\ \hline
$d^2$ & 0,3845 & 0,0260 & \textbf{0,0221} & 0,6243 & 0,8837\\ \hline
$d$ & 0,620 & 0,161 & \textbf{0,149} & 0,790 & 0,940\\ \hline
\end{tabular}
\end{center}
Minimum bei $\mathbf{s_0 = 2{,}5}$ (TPR 0,86 / FPR 0,05), knapp vor $s_0=2{,}0$.

\textbf{Vergleich mit $s_0=1{,}0$:} TPR 0,99 statt 0,86, also $0{,}13\cdot80 \approx 10$
zusätzlich gefundene Ausreißer. Fehlalarme: $0{,}05\cdot4920 \approx 246$ gegen
$0{,}62\cdot4920 \approx 3050$, also rund \textbf{2800 zusätzliche Fehlalarme}.

\neu \textbf{Auf einen Ausreißer heruntergerechnet:} $2800/10 \approx \mathbf{280}$ Fehlalarme je
zusätzlich gefundenem Ausreißer. Genau danach fragt „was kostet der letzte Rest TPR`` -- eine
\textbf{Verhältniszahl}, nicht zwei absolute Zahlen nebeneinander. Damit ist die Antwort in einem
Satz begründbar: \emph{lohnt sich nur, wenn ein übersehener Ausreißer mindestens 280-mal so teuer
ist wie ein Fehlalarm.}
\end{loes}

\textbf{e)} Warum ist das Histogramm die schlechteste der drei Darstellungen? \pkt{4}

\begin{loes}
Es zeigt nur, \emph{dass} die Masse normal ist, nicht \emph{wo} sie aufhört. Die Region, in der die
Entscheidung fällt, ist genau die, in der die Balken schon auf der Nulllinie liegen.
\begin{itemize}
\item \textbf{Sortierter Score-Plot} (Score über dem Rang): die \textbf{Übersetzungstabelle
  Schwellwert $\leftrightarrow$ Anzahl}. Schwelle auf der $y$-Achse ablesen, Indexposition auf der
  $x$-Achse, $N -$ Index $=$ Alarme. Eine \textbf{Schulter} ist meist eine zweite Normalgruppe;
  erst wo die Kurve in die \textbf{Senkrechte} geht, beginnen die Ausreißer.
\item \textbf{ROC-Kurve} (sobald Labels da sind): Betriebspunkt am \textbf{Knie}, und zu jedem
  Schwellwert das Paar (FPR, TPR), also die Kosten in Fehlalarmen. Die \textbf{AUC} daneben trennt
  „Modell schlecht`` von „Schwelle schlecht``.
\end{itemize}
(Ebenfalls gültig: der \textbf{Boxplot} nach Klassen -- er macht den schweren Oberschwanz aus c)
sichtbar.)
\end{loes}

\newpage
\section*{Aufgabe 6 -- SVM und One-Class SVM \pkt{22}}

\subsection*{6.1 a) Trenngerade bei sehr großem $C$ \pkt{4}}

\begin{center}
\begin{tikzpicture}[scale=0.72]
  \draw[step=1,gray!30,line width=0.3pt] (0,0) grid (8,6);
  \draw[->] (0,0) -- (8.4,0) node[right,scale=0.85]{$x_1$};
  \draw[->] (0,0) -- (0,6.4) node[above,scale=0.85]{$x_2$};
  \foreach \x in {1,...,8} \draw (\x,0) -- (\x,-0.12) node[below,scale=0.7]{\x};
  \foreach \y in {1,...,6} \draw (0,\y) -- (-0.12,\y) node[left,scale=0.7]{\y};
  \draw[gray,dashed,line width=0.7pt] (0,5) -- (5,0);
  \draw[gray,dashed,line width=0.7pt] (1,6) -- (7,0);
  \draw[loesfarbe,line width=1.3pt] (0.5,5.5) -- (6,0);
  \node[scale=0.75,loesfarbe] at (7.0,1.5) {$x_1+x_2=6$};
  \foreach \p in {(1,2),(2,1),(3,1)} \fill \p circle (0.11);
  \fill (2,3) circle (0.11);
  \draw[loesfarbe,line width=1pt] (2,3) circle (0.26);
  \foreach \p in {(5,4),(6,3),(6,5),(7,4)}
     \draw[line width=1pt] \p ++(-0.13,-0.13) -- ++(0.26,0.26) \p ++(-0.13,0.13) -- ++(0.26,-0.26);
  \draw[line width=1pt] (3,4) ++(-0.13,-0.13) -- ++(0.26,0.26) (3,4) ++(-0.13,0.13) -- ++(0.26,-0.26);
  \draw[loesfarbe,line width=1pt] (3,4) circle (0.26);
  \draw[line width=1pt] (4,3) ++(-0.13,-0.13) -- ++(0.26,0.26) (4,3) ++(-0.13,0.13) -- ++(0.26,-0.26);
  \node[scale=0.7] at (4.35,2.55) {$(4|3)$};
\end{tikzpicture}
\end{center}

\begin{loes}
Sehr großes $C$ $=$ \textbf{harte Margin}, kein Punkt darf in den Streifen.

\textbf{Support-Vektoren:} $\bullet(2|3)$ und $\times(3|4)$ -- das nächstgelegene Paar aus
verschiedenen Klassen (Abstand $\sqrt2$), oben eingekreist.

\textbf{Rechenweg:} $\boldsymbol{w}$ steht senkrecht auf ihrer Verbindung, also
$\boldsymbol{w}\propto(3|4)-(2|3) = (1|1)$; die Gerade geht durch den Mittelpunkt $(2{,}5|3{,}5)$:
\[
g(\boldsymbol{x}) = x_1+x_2-6, \qquad g(2|3) = -1, \quad g(3|4) = +1
\]
Die Normierung $\pm1$ stimmt damit bereits, also $\|\boldsymbol{w}\| = \sqrt2$ und
\[
\text{Margin-Breite} \;=\; \frac{2}{\|\boldsymbol{w}\|} \;=\; \frac{2}{\sqrt2} \;=\; \sqrt2
\;\approx\; 1{,}41
\]
Probe: alle $\bullet$ erfüllen $x_1+x_2\le5$ (3, 3, 5, 4), alle $\times$ erfüllen $x_1+x_2\ge7$
(7, 9, 9, 11, 11). \checkmark
\end{loes}

\subsection*{6.1 b) $C$ wird schrittweise kleiner \pkt{5}}

\begin{loes}
\begin{itemize}
\item \textbf{Margin:} wird \textbf{breiter} -- Verletzungen sind billiger, ein großer Abstand
  lohnt sich mehr als ein sauberer Streifen.
\item \textbf{Lage der Trenngeraden:} löst sich von den beiden Randpunkten und richtet sich nach
  der \textbf{Masse} der Punkte -- stabiler gegen Ausreißer, unpräziser an der Klassengrenze.
\item \textbf{Support-Vektoren:} ihre Anzahl \textbf{steigt} -- jeder Punkt \emph{in} der Margin
  und jeder Verletzer ist per Definition SV.
\item \textbf{Fehlklassifikationen:} treten auf und nehmen zu, bewusst in Kauf genommen
  (Soft Margin).
\item \textbf{Grenzfall $C\to0$:} Der Fehlerterm hat kein Gewicht mehr, minimiert wird nur noch
  $\|\boldsymbol{w}\|^2$ $\Rightarrow$ $\boldsymbol{w}\to\boldsymbol{0}$, die Margin wird unendlich
  breit, die \textbf{Lage der Geraden ist beliebig} (nicht mehr an die Daten gefesselt), fast alle
  Punkte sind SV. Es bleibt kein brauchbarer Klassifikator übrig.
\end{itemize}
\end{loes}

\subsection*{6.1 c) Neuer Punkt $(4|3)$ \pkt{3}}

\begin{loes}
Einsetzen statt schätzen: $g(4|3) = 4+3-6 = +1 > 0$ $\Rightarrow$ \textbf{Klasse $\times$
($y=+1$)}. Gleichwertig über Abstände: nächster $\times$ ist $(3|4)$ mit $\sqrt2\approx1{,}41$,
nächster $\bullet$ ist $(2|3)$ mit $2$.

\neu \textbf{Lage:} $g = +1$ heißt \textbf{exakt auf dem Margin-Rand} $h_2$ ($x_1+x_2 = 7$), also
gerade \emph{nicht} innerhalb der Margin. Wäre er ein Trainingspunkt gewesen, wäre er ein weiterer
Support-Vektor. Das ist die dritte mögliche Antwort neben „innerhalb`` und „außerhalb`` und fällt
automatisch heraus, wenn man $g$ ausrechnet:
\[
|g| < 1 \Rightarrow \text{in der Margin}, \qquad
|g| = 1 \Rightarrow \text{auf dem Rand}, \qquad
|g| > 1 \Rightarrow \text{außerhalb}
\]
\end{loes}

\subsection*{6.2 a) Kernel-Trick nachrechnen \pkt{4}}

\begin{loes}
$\varphi(\boldsymbol{x}) = (x_1^2,\ \sqrt2\,x_1x_2,\ x_2^2)$, \
$\boldsymbol{a} = (1|2)$, $\boldsymbol{b} = (3|1)$.
\[
\varphi(\boldsymbol{a}) = (1,\ 2\sqrt2,\ 4), \qquad
\varphi(\boldsymbol{b}) = (9,\ 3\sqrt2,\ 1)
\]
\[
\varphi(\boldsymbol{a})\cdot\varphi(\boldsymbol{b}) = 1\cdot9 + 2\sqrt2\cdot3\sqrt2 + 4\cdot1
= 9+12+4 = \mathbf{25}
\]
\[
k(\boldsymbol{a},\boldsymbol{b}) = (\boldsymbol{a}\cdot\boldsymbol{b})^2 = (1\cdot3+2\cdot1)^2
= 5^2 = \mathbf{25}
\]
\textbf{Vergleich:} identisch. Der Kernel liefert das Skalarprodukt der \textbf{transformierten}
Punkte, ohne die Transformation je auszuführen -- links drei Multiplikationen im 3D-Raum plus zwei
Wurzeln, rechts ein Skalarprodukt im 2D-Raum und ein Quadrat.
\end{loes}

\subsection*{6.2 b) Wo Skalarprodukte genügen; RBF vs.\ polynomial \pkt{5}}

\begin{loes}
\textbf{Die zwei Rechnungen:}
\begin{enumerate}
\item das \textbf{duale Optimierungsproblem beim Training} -- dort stehen die Trainingspunkte nur
  als Paare $\boldsymbol{x}_i\cdot\boldsymbol{x}_j$;
\item die \textbf{Entscheidungsfunktion beim Anwenden} -- dort steht der neue Punkt nur als
  $\boldsymbol{x}\cdot\boldsymbol{x}_i$ gegen die Support-Vektoren.
\end{enumerate}
$\boldsymbol{w}$ selbst kommt in beiden nicht mehr vor (es fällt über
$\partial L/\partial\boldsymbol{w} = 0$ heraus) -- genau deshalb funktioniert der Trick.

\textbf{Rechnerischer Unterschied:} Der \textbf{polynomiale} Kernel bildet erst das
\textbf{Skalarprodukt} und potenziert es dann; der \textbf{RBF} setzt den \textbf{Abstand} der
beiden Punkte in eine Gaußfunktion ein.

\textbf{Warum beim RBF die einzige Möglichkeit:} Zum polynomialen Kernel gehört eine Transformation
mit \textbf{endlich vielen} Komponenten -- man könnte $\varphi(\boldsymbol{x})$ hinschreiben, es
wäre nur teurer (genau das zeigt a)). Zum RBF gehört über die Reihenentwicklung der e-Funktion eine
Transformation mit \textbf{unendlich vielen} Komponenten; sie lässt sich gar nicht ausrechnen.
\end{loes}

\subsection*{6.3 a) Lineare OCSVM, $\nu = 0{,}05$ \pkt{3}}

\begin{center}
\begin{tikzpicture}[scale=0.62]
  \draw[step=1,gray!30,line width=0.3pt] (0,0) grid (8,7);
  \draw[->] (0,0) -- (8.4,0) node[right,scale=0.85]{$x_1$};
  \draw[->] (0,0) -- (0,7.4) node[above,scale=0.85]{$x_2$};
  \foreach \x in {1,...,8} \draw (\x,0) -- (\x,-0.12) node[below,scale=0.7]{\x};
  \foreach \y in {1,...,7} \draw (0,\y) -- (-0.12,\y) node[left,scale=0.7]{\y};
  \foreach \p in {(4,4),(4.6,5),(5,4.3),(5.5,5.2),(4.2,5.6),(5.8,4.1),(5.2,5.8),(4.8,4.8),
                  (6.2,5.0),(3.8,4.9),(5.6,3.9),(4.4,4.2),(6.0,5.7),(5.0,5.3),(4.9,6.1)}
     \fill \p circle (0.1);
  \fill[black] (2.6,6.4) circle (0.1);
  \node[scale=0.65] at (2.6,6.9) {Randpunkt};
  \draw[loesfarbe,line width=1.3pt] (1.1,7) -- (8,0.1);
  \node[scale=0.65,rotate=-45,loesfarbe] at (2.5,5.2) {Trennebene};
  \node[scale=0.72] at (1.6,2.2) {\textbf{anomal}};
  \node[scale=0.62] at (1.6,1.5) {(Ursprungsseite)};
  \node[scale=0.72] at (6.7,6.6) {\textbf{normal}};
  \draw[->,line width=0.6pt] (0.25,0.25) -- (1.7,1.7);
\end{tikzpicture}
\end{center}

\begin{loes}
Die Trennebene steht \textbf{zwischen Ursprung und Datenwolke} und wird so weit wie möglich
\textbf{vom Ursprung weg} geschoben. Anomal ist die \textbf{Ursprungsseite}.
$\nu = 0{,}05$ bei 16 Punkten heißt: höchstens etwa \textbf{ein} Punkt darf dort liegen -- die
Ebene schneidet gerade den ursprungsnächsten Punkt der Wolke ab.

\textbf{Wogegen abgegrenzt wird:} Es gibt keine zweite Klasse, also nimmt die OCSVM den
\textbf{Ursprung als Stellvertreter für alles Anomale} und maximiert den Abstand der Normaldaten
zu ihm.

\neu \textbf{Der Zusatzsatz, der die Aufgabe knackt:} Der eingezeichnete Randpunkt liegt sichtbar
abseits der Wolke, aber \textbf{nicht ursprungsnäher} als sie -- der lineare Kernel stuft ihn
deshalb als \textbf{normal} ein. Er schneidet nämlich nur \emph{eine} Richtung ab: „zu nah am
Ursprung``. Alles, was \textbf{seitlich} abseits liegt, sieht er nicht. Das ist die eigentliche
Schwäche, schärfer als „er kann nur Geraden``.
\end{loes}

\subsection*{6.3 b) Größeres $\nu$; Over-/Underfitting \pkt{3}}

\begin{loes}
Die Ebene wandert \textbf{weiter vom Ursprung weg, in die Wolke hinein}: mehr Punkte landen auf der
Ursprungsseite und gelten als anomal, die Zahl der Support-Vektoren steigt (untere Schranke!), die
beschriebene Normalregion schrumpft.
\begin{itemize}
\item $\nu$ \textbf{zu klein} $\Rightarrow$ die Grenze umschließt restlos alles, auch Randpunkte
  und Rauschen $\Rightarrow$ \textbf{Overfitting}: echte Anomalien rutschen später durch.
\item $\nu$ \textbf{zu groß} $\Rightarrow$ zu viele Normaldaten abgeschnitten $\Rightarrow$
  \textbf{Underfitting} und hohe Fehlalarmrate.
\end{itemize}
\end{loes}

\subsection*{6.3 c) Wechsel auf RBF \pkt{4}}

\begin{loes}
\textbf{Entscheidungsgrenze:} keine Gerade mehr, sondern eine \textbf{geschlossene, krumme Kurve um
die Daten}. Bei mehreren getrennten Gruppen zerfällt sie in \textbf{mehrere Inseln} -- die
eigentliche Stärke gegenüber einer einzelnen Ellipse.

\textbf{$\gamma$ sehr groß:} Die Glocke wird so schmal, dass jeder Trainingspunkt nur noch
\textbf{sich selbst} erkennt; die Grenze zerfällt in winzige Kreise um die einzelnen
Trainingspunkte. Das ist die Overfitting-Signatur von $\gamma$.

\textbf{Balanced accuracy $= 0{,}5$:} Jeder neue Punkt fällt aus jedem dieser Kreise heraus und wird
als Anomalie geflaggt. Damit TPR $=1$, TNR $=0$, also $\tfrac12(1+0) = 0{,}5$ -- exakt der Wert für
\textbf{Raten}.
\end{loes}

\newpage
\section*{Aufgabe 7 -- Verfahrenswahl bei mehreren Clustern \pkt{12}}

\textbf{a)} Zwei geeignete, ein ungeeignetes Verfahren -- jedes aus der Form der Daten begründet. \pkt{6}

\begin{loes}
Datenlage: die Normaldaten bilden \textbf{drei getrennte, kompakte Gruppen} (multimodal), dazwischen
große leere Bereiche. Novelty-Setup, also Fit auf sauberen Normaldaten.

\textbf{Geeignet 1 -- OCSVM mit RBF-Kernel.} Die Entscheidungsgrenze ist eine geschlossene krumme
Kurve, die bei getrennten Gruppen in \textbf{drei Inseln} zerfällt; der leere Bereich dazwischen
bleibt außerhalb der Normalregion.

\textbf{Geeignet 2 -- abstandsbasiert (kNN-Abstand).} Der Score ist der mittlere Abstand zu den $k$
nächsten Nachbarn -- \textbf{keine Formannahme}, beliebig viele Gruppen darstellbar. Normal ist,
wer nah an \emph{irgendeiner} Gruppe liegt. (Gleichwertig: \textbf{KDE} oder \textbf{GMM mit 3
Komponenten}.)

\textbf{Ungeeignet -- Elliptic Envelope / Mahalanobis.} Es wird \textbf{eine einzige}
Normalverteilung geschätzt, die Normalregion ist \textbf{eine Ellipse}. Sie legt sich über alle
drei Gruppen und \textbf{überdeckt den leeren Bereich in der Mitte mit}. Die Formannahme
(unimodal, elliptisch) ist hier schlicht falsch.
\end{loes}

\textbf{b)} Bewertung von $P_1$ und $P_2$; welcher Punkt trennt? \pkt{4}

\begin{loes}
\begin{center}\small
\begin{tabular}{|l|c|c|}
\hline
& $P_1$ (Mitte, zwischen den Gruppen) & $P_2$ (weit außen rechts oben)\\ \hline
OCSVM (RBF) & \textbf{anomal} & anomal\\ \hline
kNN-Abstand & \textbf{anomal} & anomal\\ \hline
Elliptic Envelope & \textbf{normal} (!) & anomal\\ \hline
\end{tabular}
\end{center}
\textbf{Der trennende Punkt ist $P_1$.} Er liegt im \textbf{Loch} der multimodalen Verteilung: weit
von jeder Gruppe, aber \textbf{nahe am gemeinsamen Mittelwert}.

\neu \textbf{Die Zuspitzung, die den Punkt holt:} Die Elliptic Envelope übersieht $P_1$ nicht
einfach -- sie misst den Mahalanobis-Abstand zum Gesamtmittelwert, und der ist dort
\textbf{minimal}. Sie hält $P_1$ also für den \textbf{normalsten Punkt des ganzen Datensatzes}.
Das ist nicht „knapp daneben``, sondern die maximale Fehlbewertung.

$P_2$ trennt nichts: er liegt außerhalb jeder Beschreibung, auch außerhalb der großen Ellipse, und
wird von allen drei Verfahren gefangen.
\end{loes}

\textbf{c)} Welcher Metaparameter entscheidet über drei Inseln vs.\ eine Region? \pkt{4}

\begin{loes}
Beim ersten geeigneten Verfahren (OCSVM-RBF) ist es \textbf{$\gamma$} -- die Breite der Glocke.
\begin{itemize}
\item \textbf{$\gamma$ zu klein} (Glocken zu breit): die drei Gruppen \textbf{verschmelzen zu einer
  großen Region}, der leere Bereich wird mit abgedeckt $\Rightarrow$ Underfitting, $P_1$ rutscht
  als normal durch -- das Verfahren verhält sich dann wie die Ellipse.
\item \textbf{$\gamma$ zu groß} (Glocken zu schmal): die Regionen zerfallen in immer kleinere
  Inseln, im Extremfall ein Kreis um jeden Trainingspunkt $\Rightarrow$ Overfitting, jeder neue
  Punkt wird Anomalie, balanced accuracy 0,5.
\end{itemize}
(Bei kNN übernimmt \textbf{$k$} diese Rolle: zu groß mittelt über eine ganze Gruppe hinweg,
$k=1$ macht das Modell rauschempfindlich.)
\end{loes}

\textbf{d)} 200 Merkmale, $10^6$ Punkte -- welches Verfahren, welche zwei Gegengründe? \pkt{3}

\begin{loes}
Wahl: \textbf{Isolation Forest}.
\begin{enumerate}
\item \textbf{Skalierung.} Der Trainingsaufwand der OCSVM wächst etwa quadratisch bis kubisch mit
  der Punktzahl, und die Entscheidungsfunktion summiert bei jeder Abfrage über \emph{alle}
  Support-Vektoren -- bei $10^6$ Punkten unbrauchbar. (kNN ist im Training billig, in der Anwendung
  genauso langsam.)
\item \textbf{Curse of Dimensionality.} Bei 200 Merkmalen werden alle paarweisen Abstände einander
  ähnlich; RBF-Kernel und kNN-Abstand verlieren ihre Trennschärfe.
\end{enumerate}
Der Isolation Forest benutzt \textbf{keine Abstände}, sondern achsenparallele Splits je ein Merkmal,
arbeitet mit Subsampling ($n_{\text{sub}} = 256$) und skaliert linear in $N$. Er braucht auch keine
Normalisierung.
\end{loes}

\newpage
\section*{Aufgabe 8 -- Kurzfragen quer durch die Vorlesung \pkt{10}}

\textbf{a)} AnoGAN-Score, warum eine Optimierung nötig ist, was f-AnoGAN löst. \pkt{5}

\begin{loes}
\textbf{Berechnung.} Ein GAN kann nur in eine Richtung: $z \to$ Bild. Für ein neues Bild $x$ ist
der Latentvektor unbekannt. Man hält deshalb \textbf{G und D fest} und sucht per
\textbf{Gradientenabstieg im Latentraum} (im Praktikum $\approx500$ Iterationen) das $z$, dessen
erzeugtes Bild $G(z)$ dem $x$ am ähnlichsten ist. Der Score ist die \textbf{verbleibende
Restabweichung} aus zwei Teilen: dem \textbf{Residual-Anteil} (Unterschied zwischen $x$ und $G(z)$
im Bild) und dem \textbf{Discriminator-Anteil} (Unterschied der Merkmale aus einer mittleren
D-Schicht, \texttt{layers[-2]}), gewichtet mit $\lambda = 0{,}1$. Normale Bilder kann G gut
nachbauen $\Rightarrow$ kleiner Rest; Anomalien liegen außerhalb der gelernten Verteilung
$\Rightarrow$ großer Rest.

\textbf{Warum eine Optimierung nötig ist.} Das GAN besitzt \textbf{kein inverses Mapping}
Bild $\to z$ -- der Encoder-Weg fehlt konstruktionsbedingt. Das $z$ muss deshalb für \emph{jedes
einzelne Testbild neu gesucht} werden; genau das macht AnoGAN in der Anwendung unpraktikabel
langsam.

\textbf{Was f-AnoGAN löst.} Es trainiert \textbf{zusätzlich einen Encoder} $E$, der $x$ direkt auf
$z$ abbildet (izif-Training, G und D dabei fest). Danach ist der Score in \textbf{einem einzigen
Vorwärtsdurchlauf} $x \to E(x) \to G(E(x))$ berechenbar -- aus hunderten Optimierungsschritten je
Bild wird ein Netzaufruf.
\end{loes}

\textbf{b)} Zwei Diversitätsquellen des RandNet-Ensembles; warum Median statt Mittelwert? \pkt{4}

\begin{loes}
\textbf{Quelle 1 -- zufällige Verbindungsmasken.} Jede Schicht jedes Netzes bekommt
\textbf{einmalig und fest} eine eigene binäre Maske auf die Gewichtsmatrix ($\approx37\,\%$ der
Verbindungen gekappt, $\approx63\,\%$ überleben) $\Rightarrow$ jedes Netz hat eine \textbf{andere
Architektur}. \emph{Nicht Dropout:} das trifft \textbf{Neuronen}, wird \textbf{pro Schritt neu}
gewürfelt und ist in der Inferenz aus.

\textbf{Quelle 2 -- unterschiedliche Trainingsdaten.} Jedes Netz trainiert auf einer anderen
zufälligen \textbf{Teilmenge} (je etwa 10\,\% der Daten, Bagging).
(Schwächer hinzu: die zufällige Gewichtsinitialisierung.)

\textbf{Median statt Mittelwert.} Die Rekonstruktionsfehler werden je Netz auf Standardabweichung 1
normiert und dann aggregiert. Ein einzelnes schlecht konvergiertes Netz kann für ein Sample einen
\textbf{extremen} Fehler liefern; der Mittelwert würde davon nach oben gezogen, der \textbf{Median
ist robust} dagegen. Genau diese Robustheit \emph{ist} die versprochene Stabilisierung des Scores.
\end{loes}

\textbf{c)} Vergleich GOAD und Deep SVDD.

\begin{center}\small
\begin{tabular}{|p{3.2cm}|p{5.7cm}|p{5.7cm}|}
\hline
 & \textbf{Deep SVDD} & \textbf{GOAD}\\ \hline
Was bilden die Normaldaten im Latent Space?
& \textcolor{loesfarbe}{\textbf{Eine} kompakte Kugel um ein festes Zentrum $c$.}
& \textcolor{loesfarbe}{\textbf{$M$ Cluster} -- einen je Transformation, jeder um sein eigenes
Zentrum.}\\ \hline
Woher kommt das Lernsignal?
& \textcolor{loesfarbe}{Allein aus dem \textbf{Abstand zu $c$}; die Zielfunktion zieht alle
Normaldaten an $c$ heran. Keine Labels.}
& \textcolor{loesfarbe}{Aus \textbf{selbst erzeugten Pseudo-Labels}: welche der $M$
Transformationen wurde angewendet? Kreuzentropie darauf $+$ Triplet Center Loss $+$ $L_2$-Term auf
die Latent-Vektoren.}\\ \hline
Was ist die triviale Lösung, und was verhindert sie?
& \textcolor{loesfarbe}{Die \textbf{konstante Abbildung} (Netz auf 0). Verhindert durch: kein Bias,
keine gedeckelten Aktivierungen (LeakyReLU statt sigmoid/tanh), $c \neq 0$ und \textbf{einmal vor
dem Training} gesetzt, danach nie geändert.}
& \textcolor{loesfarbe}{Alle Transformationen auf \textbf{denselben} Punkt abzubilden. Verhindert
durch die \textbf{Kreuzentropie}: die $M$ Cluster müssen unterscheidbar bleiben, während $L_{TC}$
sie gleichzeitig dicht zusammenzieht.}\\ \hline
Wie entsteht der Anomalie-Score?
& \textcolor{loesfarbe}{\textbf{Abstand von $\varphi(x)$ zu $c$}. Der Radius $R$ als
$(1-\nu)$-Quantil der Trainingsabstände legt die $\pm1$-Grenze fest -- verschiebt nur den
Betriebspunkt, nicht die AUC.}
& \textcolor{loesfarbe}{Alle $M$ Transformationen auf $x$ anwenden, jede Version im Latent Space
verorten und prüfen, ob sie bei \textbf{ihrem eigenen} Zentrum landet; Score $=$ Summe der
negativen Log-Wahrscheinlichkeiten.}\\ \hline
\end{tabular}
\end{center}

\vspace{6pt}
\noindent\rule{\linewidth}{0.8pt}
\begin{center}\small
Selbst erstellte Lösungen zur generierten Probeklausur 2 -- keine offizielle Musterlösung.
\end{center}

\end{document}
